% Generated by the blueprint tool (blueprint-pdf/1) from a blueprint snapshot; edit the
% Markdown sources, not this file.  Compile with `latexmk -pdf` or by running
% `pdflatex` twice (the second run settles the contents and the page references).
% Images, if any, are read from assets/ next to this file.
\documentclass[11pt]{article}
\usepackage[T1]{fontenc}
\usepackage[utf8]{inputenc}
\usepackage{lmodern}
\usepackage{textcomp}
\usepackage[a4paper,margin=2.5cm,headheight=14pt]{geometry}
\usepackage{amsmath,amssymb,amsthm}
\usepackage{amscd}  % KaTeX's CD diagrams; part of amsmath, so always there
\usepackage{mathtools}
\usepackage{bm}
\usepackage{graphicx}
\usepackage{xcolor}
\usepackage{array,booktabs,tabularx}
\usepackage{fancyvrb}
\usepackage[normalem]{ulem}
\usepackage{microtype}
% Optional, and only with their fonts: a small TeX installation has the
% packages without the fonts, which fails at the first formula.
\IfFileExists{mathrsfs.sty}{\IfFileExists{rsfs10.tfm}{\usepackage{mathrsfs}}{}}{}
\IfFileExists{stmaryrd.sty}{\IfFileExists{stmary10.tfm}{\usepackage{stmaryrd}}{}}{}
\IfFileExists{cancel.sty}{\usepackage{cancel}}{}
\usepackage[hyphens]{url}
\usepackage[pdfusetitle,bookmarksnumbered=false]{hyperref}
\usepackage{bookmark}
\definecolor{bplink}{HTML}{1F4E9C}
\definecolor{bpmuted}{HTML}{5F676E}
\hypersetup{colorlinks=true,linkcolor=bplink,urlcolor=bplink,citecolor=bplink}
\setlength{\parindent}{0pt}
\setlength{\parskip}{0.55em plus 0.2em minus 0.1em}
% Lists as tight as the page's.
\makeatletter
\def\@listi{\leftmargin\leftmargini\parsep 0pt\topsep 0.3em plus 0.1em\itemsep 0.1em plus 0.1em}
\let\@listI\@listi
\def\@listii{\leftmargin\leftmarginii\labelwidth\leftmarginii\advance\labelwidth-\labelsep\parsep 0pt\topsep 0.1em\itemsep 0.1em}
% An ordered list that does not start at 1.
\newcommand{\bpstartat}[1]{\setcounter{\@enumctr}{\numexpr#1-1\relax}}
\makeatother
\setcounter{tocdepth}{3}
\emergencystretch=3em
\allowdisplaybreaks

% A numbered statement, set as a paper sets one: "Theorem 2.1 (Title)."
% #1 is the lead ("Theorem~2.1"), #2 what follows it in roman.
\newenvironment{bpstatement}[2]{\par\addvspace{0.6\baselineskip}\noindent\textbf{#1}#2\textbf{.}\enspace\ignorespaces}
  {\par\addvspace{0.4\baselineskip}}
\newenvironment{bpstatementit}[2]{\begin{bpstatement}{#1}{#2}\itshape\ignorespaces}{\end{bpstatement}}
% A step: a relation with prose of its own, read under the object it starts from.
\newenvironment{bpstep}[1]{\par\addvspace{0.3\baselineskip}\begin{list}{}{\leftmargin=1.6em\rightmargin=0pt\topsep=0pt}\item\relax
  {\small\color{bpmuted}#1}\par\nopagebreak}{\end{list}\par}
% Notes about where an object is written out.
\newenvironment{bpnote}{\par\small\color{bpmuted}\itshape}{\par}
\newcommand{\bpomitted}[1]{\par{\small\color{bpmuted}\itshape #1\par}}
\newcommand{\bpbodyheading}[1]{\par\addvspace{0.4\baselineskip}\noindent\textbf{#1}\par\nopagebreak}
\newcommand{\bprule}{\par\medskip\noindent\hfil\rule{0.3\linewidth}{0.4pt}\hfil\par\medskip}
\newcommand{\bpboundary}[1]{\par\nopagebreak{\small\color{bpmuted}#1\par}\nopagebreak}
% A reference to something this export leaves out, and one that resolves nowhere.
\newcommand{\bpout}[1]{#1\textsuperscript{\hyperlink{bp-outside}{\ensuremath{\ast}}}}
\newcommand{\bpbroken}[1]{\textcolor{red!70!black}{[#1?]}}
\newcommand{\bpslug}[1]{\texttt{#1}}
% A bibliography entry: label, anchor, entry; the label hangs in the margin.
\newcommand{\bpbibitem}[3]{\par\setbox0\hbox{[#1]\enspace}\ifdim\wd0<5em\setbox0\hbox to 5em{[#1]\hfil}\fi
  \hangindent=\wd0 \hangafter=1 \noindent\box0\phantomsection\label{#2}#3\par\smallskip}
% A character the export knows no way to set.
\newcommand{\bpglyph}[1]{\fbox{\scriptsize #1}}
% Images, at their own size unless they are wider than the line.
\newsavebox{\bpimagebox}
\newcommand{\bpimage}[1]{\sbox{\bpimagebox}{\includegraphics{#1}}%
  \ifdim\wd\bpimagebox>\linewidth\resizebox{\linewidth}{!}{\usebox{\bpimagebox}}\else\usebox{\bpimagebox}\fi}
\newcommand{\bpmissingimage}[1]{\fbox{\parbox{0.8\linewidth}{\centering\small\color{bpmuted}Image not included: #1}}}
\DefineVerbatimEnvironment{bpverbatim}{Verbatim}{fontsize=\small,xleftmargin=1em}
\renewcommand{\qedsymbol}{\ensuremath{\blacksquare}}
\hypersetup{pdftitle={Weil Conjectures — A Nagata compactification and Chow's lemma},pdfsubject={Appendix A Nagata compactification and Chow's lemma; all levels},pdfcreator={blueprint}}
% KaTeX commands the website knows
\providecommand{\R}{\mathbb{R}}\providecommand{\N}{\mathbb{N}}\providecommand{\Z}{\mathbb{Z}}
\providecommand{\Q}{\mathbb{Q}}\providecommand{\C}{\mathbb{C}}
\providecommand{\Reals}{\mathbb{R}}\providecommand{\reals}{\mathbb{R}}\providecommand{\Complex}{\mathbb{C}}
\providecommand{\cnums}{\mathbb{C}}\providecommand{\natnums}{\mathbb{N}}\providecommand{\Bbb}[1]{\mathbb{#1}}
\providecommand{\bold}[1]{\mathbf{#1}}\providecommand{\frak}[1]{\mathfrak{#1}}\providecommand{\mathbbm}[1]{\mathbb{#1}}
\providecommand{\lang}{\langle}\providecommand{\rang}{\rangle}\providecommand{\bra}[1]{\langle #1\rvert}
\providecommand{\ket}[1]{\lvert #1\rangle}\providecommand{\braket}[1]{\langle #1\rangle}
\providecommand{\set}[1]{\{#1\}}\providecommand{\Set}[1]{\left\{#1\right\}}
\providecommand{\infin}{\infty}\providecommand{\isin}{\in}\providecommand{\plusmn}{\pm}\providecommand{\sub}{\subset}
\providecommand{\sube}{\subseteq}\providecommand{\supe}{\supseteq}\providecommand{\exist}{\exists}
\providecommand{\bull}{\bullet}\providecommand{\larr}{\leftarrow}\providecommand{\rarr}{\rightarrow}
\providecommand{\lrarr}{\leftrightarrow}\providecommand{\Larr}{\Leftarrow}\providecommand{\Rarr}{\Rightarrow}
\providecommand{\Lrarr}{\Leftrightarrow}\providecommand{\uarr}{\uparrow}\providecommand{\darr}{\downarrow}
\providecommand{\harr}{\leftrightarrow}\providecommand{\lArr}{\Leftarrow}\providecommand{\rArr}{\Rightarrow}
\providecommand{\hArr}{\Leftrightarrow}\providecommand{\alef}{\aleph}\providecommand{\alefsym}{\aleph}
\providecommand{\real}{\Re}\providecommand{\image}{\Im}\providecommand{\weierp}{\wp}
\providecommand{\Alpha}{\mathrm{A}}\providecommand{\Beta}{\mathrm{B}}\providecommand{\Epsilon}{\mathrm{E}}
\providecommand{\Zeta}{\mathrm{Z}}\providecommand{\Eta}{\mathrm{H}}\providecommand{\Iota}{\mathrm{I}}
\providecommand{\Kappa}{\mathrm{K}}\providecommand{\Mu}{\mathrm{M}}\providecommand{\Nu}{\mathrm{N}}
\providecommand{\Omicron}{\mathrm{O}}\providecommand{\Rho}{\mathrm{P}}\providecommand{\Tau}{\mathrm{T}}
\providecommand{\Chi}{\mathrm{X}}\providecommand{\omicron}{o}\providecommand{\gt}{>}\providecommand{\lt}{<}
\providecommand{\And}{\&}\providecommand{\hdashline}{\hline}\providecommand{\operatornamewithlimits}[1]{\operatorname*{#1}}
\providecommand{\llbracket}{[\![}\providecommand{\rrbracket}{]\!]}\providecommand{\cancel}[1]{#1}
\providecommand{\mathscr}{\mathcal}
% the characters beyond ASCII this document uses
\DeclareUnicodeCharacter{2160}{\textup{I}}
% the project's macros ([katex.macros] in blueprint.toml)
\providecommand{\A}{}\renewcommand{\A}{\mathbf A}
\providecommand{\AEn}{}\renewcommand{\AEn}{A_{E,n}}
\providecommand{\Aaf}{}\renewcommand{\Aaf}{\mathbf A}
\providecommand{\Ab}{}\renewcommand{\Ab}{\mathsf{Ab}}
\providecommand{\AffSch}{}\renewcommand{\AffSch}{\operatorname{AffSch}}
\providecommand{\Aone}{}\renewcommand{\Aone}{\mathbf A^1}
\providecommand{\Bl}{}\renewcommand{\Bl}{\operatorname{Bl}}
\providecommand{\Br}{}\renewcommand{\Br}{\mathrm{Br}}
\providecommand{\C}{}\renewcommand{\C}{\mathbb C}
\providecommand{\Cech}{}\renewcommand{\Cech}{\check{H}}
\providecommand{\Coext}{}\renewcommand{\Coext}{\operatorname{Coext}}
\providecommand{\Cone}{}\renewcommand{\Cone}{\operatorname{Cone}}
\providecommand{\Cons}{}\renewcommand{\Cons}{\operatorname{Cons}}
\providecommand{\Cref}{}\renewcommand{\Cref}[1]{\text{#1}}
\providecommand{\D}{}\renewcommand{\D}{D}
\providecommand{\DLAFT}{}\renewcommand{\DLAFT}{\operatorname{FT}}
\providecommand{\DLAFr}{}\renewcommand{\DLAFr}{\operatorname{Fr}}
\providecommand{\DLAQ}{}\renewcommand{\DLAQ}{{\overline{\mathbf Q}_\ell}}
\providecommand{\DLRQ}{}\renewcommand{\DLRQ}{{\overline{\mathbf Q}_\ell}}
\providecommand{\Dbc}{}\renewcommand{\Dbc}{D^b_c}
\providecommand{\Dc}{}\renewcommand{\Dc}{D_{\mathrm{cons}}}
\providecommand{\Dctf}{}\renewcommand{\Dctf}{D^b_{ctf}}
\providecommand{\Dic}{}\renewcommand{\Dic}{D_{\IndCons}}
\providecommand{\Dplus}{}\renewcommand{\Dplus}{D^{+}}
\providecommand{\End}{}\renewcommand{\End}{\operatorname{End}}
\providecommand{\Et}{}\renewcommand{\Et}{\mathrm{\acute{e}t}}
\providecommand{\Ex}{}\renewcommand{\Ex}{\mathrm{Ex}}
\providecommand{\Ext}{}\renewcommand{\Ext}{\operatorname{Ext}}
\providecommand{\F}{}\renewcommand{\F}{\mathbb F}
\providecommand{\FT}{}\renewcommand{\FT}{\operatorname{FT}}
\providecommand{\Fbar}{}\renewcommand{\Fbar}{{\overline{\mathbf F}_q}}
\providecommand{\Fell}{}\renewcommand{\Fell}{\mathbf F_\ell}
\providecommand{\Fib}{}\renewcommand{\Fib}{\operatorname{Fib}}
\providecommand{\Fix}{}\renewcommand{\Fix}{\operatorname{Fix}}
\providecommand{\Fp}{}\renewcommand{\Fp}{{\mathbf F_{p}}}
\providecommand{\Fq}{}\renewcommand{\Fq}{{\mathbf F_q}}
\providecommand{\Fqn}{}\renewcommand{\Fqn}{{\mathbf F_{q^{n}}}}
\providecommand{\Frob}{}\renewcommand{\Frob}{\operatorname{Frob}}
\providecommand{\GL}{}\renewcommand{\GL}{\operatorname{GL}}
\providecommand{\Gal}{}\renewcommand{\Gal}{\operatorname{Gal}}
\providecommand{\Hens}{}\renewcommand{\Hens}{\operatorname{Hens}}
\providecommand{\Hilb}{}\renewcommand{\Hilb}{\operatorname{Hilb}}
\providecommand{\Hom}{}\renewcommand{\Hom}{\operatorname{Hom}}
\providecommand{\IndCons}{}\renewcommand{\IndCons}{\operatorname{IndCons}}
\providecommand{\Kar}{}\renewcommand{\Kar}{\operatorname{Kar}}
\providecommand{\LT}{}\renewcommand{\LT}{\operatorname{LT}}
\providecommand{\Lie}{}\renewcommand{\Lie}{\operatorname{Lie}}
\providecommand{\Loc}{}\renewcommand{\Loc}{\operatorname{Loc}}
\providecommand{\Mod}{}\renewcommand{\Mod}{\mathsf{Mod}}
\providecommand{\N}{}\renewcommand{\N}{\mathbb{N}}
\providecommand{\Needspace}{}\renewcommand{\Needspace}{\par}
\providecommand{\O}{}\renewcommand{\O}{{\mathcal{O}}}
\providecommand{\Ob}{}\renewcommand{\Ob}{\mathrm{Ob}}
\providecommand{\Oe}{}\renewcommand{\Oe}{\mathcal O_E}
\providecommand{\PP}{}\renewcommand{\PP}{\mathbf{P}}
\providecommand{\Perf}{}\renewcommand{\Perf}{\operatorname{Perf}}
\providecommand{\Pic}{}\renewcommand{\Pic}{\mathrm{Pic}}
\providecommand{\Pm}{}\renewcommand{\Pm}{\mathbf P^{m}}
\providecommand{\Pone}{}\renewcommand{\Pone}{\mathbf P^{1}}
\providecommand{\Pqn}{}\renewcommand{\Pqn}{\mathbf P^{n}}
\providecommand{\ProFin}{}\renewcommand{\ProFin}{\operatorname{ProFin}}
\providecommand{\Proet}{}\renewcommand{\Proet}{\mathrm{pro\acute{e}t}}
\providecommand{\Q}{}\renewcommand{\Q}{\mathbb Q}
\providecommand{\Qb}{}\renewcommand{\Qb}{\overline{\mathbf Q}_{\ell}}
\providecommand{\Qbar}{}\renewcommand{\Qbar}{\overline{\mathbf Q}}
\providecommand{\QbarX}{}\renewcommand{\QbarX}{\overline{\mathbf Q}_{\ell,X}}
\providecommand{\Qbarl}{}\renewcommand{\Qbarl}{{\overline{\mathbf Q}_{\ell}}}
\providecommand{\Qell}{}\renewcommand{\Qell}{{\mathbf{Q}_{\ell}}}
\providecommand{\Ql}{}\renewcommand{\Ql}{\mathbf Q_{\ell}}
\providecommand{\Qlb}{}\renewcommand{\Qlb}{{\overline{\mathbf Q}_\ell}}
\providecommand{\Qlbar}{}\renewcommand{\Qlbar}{\overline{\mathbf Q}_\ell}
\providecommand{\R}{}\renewcommand{\R}{\mathbb{R}}
\providecommand{\RG}{}\renewcommand{\RG}{\operatorname{\mathrm{R}\Gamma}}
\providecommand{\RGamma}{}\renewcommand{\RGamma}{R\Gamma}
\providecommand{\RGammac}{}\renewcommand{\RGammac}{R\Gamma_c}
\providecommand{\RHom}{}\renewcommand{\RHom}{R\mathcal{H}om}
\providecommand{\RHomX}{}\renewcommand{\RHomX}{R\mathcal H om_X}
\providecommand{\Rat}{}\renewcommand{\Rat}{\operatorname{IsRat}}
\providecommand{\Res}{}\renewcommand{\Res}{\operatorname{Res}}
\providecommand{\Rlim}{}\renewcommand{\Rlim}{R\!\varprojlim}
\providecommand{\SIm}{}\renewcommand{\SIm}{\operatorname{SIm}}
\providecommand{\Sch}{}\renewcommand{\Sch}{\mathsf{Sch}}
\providecommand{\Sh}{}\renewcommand{\Sh}{\operatorname{Sh}}
\providecommand{\Span}{}\renewcommand{\Span}{\operatorname{span}}
\providecommand{\Spec}{}\renewcommand{\Spec}{\operatorname{Spec}}
\providecommand{\Tr}{}\renewcommand{\Tr}{\operatorname{Tr}}
\providecommand{\Trsm}{}\renewcommand{\Trsm}{\operatorname{Tr}^{\mathrm{sm}}}
\providecommand{\Wt}{}\renewcommand{\Wt}{\operatorname{Wt}}
\providecommand{\Z}{}\renewcommand{\Z}{\mathbb{Z}}
\providecommand{\Zar}{}\renewcommand{\Zar}{\mathrm{zar}}
\providecommand{\Zell}{}\renewcommand{\Zell}{{\mathbf{Z}_{\ell}}}
\providecommand{\Zl}{}\renewcommand{\Zl}{{\mathbf Z_\ell}}
\providecommand{\affet}{}\renewcommand{\affet}[1]{{#1}^{\mathrm{aff}}_{\mathrm{ét}}}
\providecommand{\affproet}{}\renewcommand{\affproet}[1]{{#1}^{\mathrm{aff}}_{\mathrm{proét}}}
\providecommand{\boxt}{}\renewcommand{\boxt}{\boxtimes_A^{\mathrm L}}
\providecommand{\cF}{}\renewcommand{\cF}{\mathcal{F}}
\providecommand{\cHom}{}\renewcommand{\cHom}{R\mathcal{H}\!om}
\providecommand{\cI}{}\renewcommand{\cI}{\mathcal{I}}
\providecommand{\cL}{}\renewcommand{\cL}{\mathcal L}
\providecommand{\cM}{}\renewcommand{\cM}{\mathcal M}
\providecommand{\cN}{}\renewcommand{\cN}{\mathcal N}
\providecommand{\cO}{}\renewcommand{\cO}{\mathcal O}
\providecommand{\cT}{}\renewcommand{\cT}{\mathcal T}
\providecommand{\checked}{}\renewcommand{\checked}{\color{red}\aftergroup\restorecolor}
\providecommand{\coeff}{}\renewcommand{\coeff}{\operatorname{coeff}}
\providecommand{\coker}{}\renewcommand{\coker}{\operatorname{coker}}
\providecommand{\colim}{}\renewcommand{\colim}{\operatorname{colim}}
\providecommand{\constructedby}{}\renewcommand{\constructedby}[1]{\par\smallskip\noindent\textcolor{purple!60!black}{\textsf{Constructed by:} }
  \footnotesize\sloppy\nolinkurl{#1}\par}
\providecommand{\deepresult}{}\renewcommand{\deepresult}[1]{\par\smallskip\noindent\textcolor{orange!45!black}{
  \textsf{Blueprint role: named deep theorem node; see the cluster proof architecture.}}
  \par\noindent\textcolor{orange!45!black}{\textsf{Mathematical source:} }#1\par}
\providecommand{\dt}{}\renewcommand{\dt}{\mathbin{\otimes^{\mathbf L}}}
\providecommand{\dw}{}\renewcommand{\dw}{\operatorname{dw}}
\providecommand{\et}{}\renewcommand{\et}[1]{{#1}_{\mathrm{ét}}}
\providecommand{\ev}{}\renewcommand{\ev}{\operatorname{ev}}
\providecommand{\fib}{}\renewcommand{\fib}{\operatorname{fib}}
\providecommand{\formalizationnote}{}\renewcommand{\formalizationnote}[1]{\begin{quote}\small\raggedright\textbf{Formalization note.} #1\end{quote}}
\providecommand{\fpqc}{}\renewcommand{\fpqc}{X_{\mathrm{fpqc}}}
\providecommand{\gammacont}{}\renewcommand{\gammacont}{\Gamma_{\mathrm{cont}}}
\providecommand{\gm}{}\renewcommand{\gm}{\mathbf G_m}
\providecommand{\hocolim}{}\renewcommand{\hocolim}{\operatorname{hocolim}}
\providecommand{\holim}{}\renewcommand{\holim}{\operatorname{holim}}
\providecommand{\hrB}{}\renewcommand{\hrB}{\@ifnextchar]{\hrEnd}{\hrStep}}
\providecommand{\hrEnd}{}\renewcommand{\hrEnd}{\@firstoftwo{}}
\providecommand{\hrStep}{}\renewcommand{\hrStep}{\@firstoftwo{\hrB}}
\providecommand{\id}{}\renewcommand{\id}{\operatorname{id}}
\providecommand{\kunFl}{}\renewcommand{\kunFl}{\mathbf F_\ell}
\providecommand{\kunZn}{}\renewcommand{\kunZn}{\mathbf Z/n\mathbf Z}
\providecommand{\kuncd}{}\renewcommand{\kuncd}{\operatorname{cd}}
\providecommand{\kundim}{}\renewcommand{\kundim}{\dim}
\providecommand{\kunotimes}{}\renewcommand{\kunotimes}[1]{\otimes^{\mathbf L}_{#1}}
\providecommand{\kunsupp}{}\renewcommand{\kunsupp}{\operatorname{Supp}}
\providecommand{\lotimes}{}\renewcommand{\lotimes}{\otimes^{\mathrm L}}
\providecommand{\m}{}\renewcommand{\m}{{\mathfrak{m}}}
\providecommand{\n}{}\renewcommand{\n}{{\mathfrak{n}}}
\providecommand{\nagataGap}{}\renewcommand{\nagataGap}[1]{\par\noindent\textbf{Formalization gap.} #1\par}
\providecommand{\openobligation}{}\renewcommand{\openobligation}[1]{\par\smallskip\noindent\textcolor{red!65!black}{\textsf{Open obligation:} }
  \footnotesize\sloppy\nolinkurl{#1}\par}
\providecommand{\p}{}\renewcommand{\p}{{\mathfrak{p}}}
\providecommand{\pr}{}\renewcommand{\pr}{\operatorname{pr}}
\providecommand{\preshv}{}\renewcommand{\preshv}[1]{\mathrm{PreShv}(#1)}
\providecommand{\pro}{}\renewcommand{\pro}[1]{\mathrm{Pro}(#1)}
\providecommand{\proet}{}\renewcommand{\proet}{X_{\mathrm{proét}}}
\providecommand{\q}{}\renewcommand{\q}{{\mathfrak{q}}}
\providecommand{\rad}{}\renewcommand{\rad}{\operatorname{rad}}
\providecommand{\rank}{}\renewcommand{\rank}{\operatorname{rank}}
\providecommand{\red}{}\renewcommand{\red}{\color{red}}
\providecommand{\restorecolor}{}\renewcommand{\restorecolor}{\color{black}}
\providecommand{\retiredobligation}{}\renewcommand{\retiredobligation}[1]{\par\smallskip\noindent\textcolor{red!75!black}{\textsf{Retired invalid route:} }
  \footnotesize\sloppy\nolinkurl{#1}\par}
\providecommand{\rk}{}\renewcommand{\rk}{\operatorname{rk}}
\providecommand{\sep}{}\renewcommand{\sep}{\textnormal{sep}}
\providecommand{\shv}{}\renewcommand{\shv}[1]{\mathrm{Shv}(#1)}
\providecommand{\spec}{}\renewcommand{\spec}[1]{\mathrm{Spec}(#1)}
\providecommand{\stacks}{}\renewcommand{\stacks}[1]{\href{https://stacks.math.columbia.edu/tag/#1}{\textnormal{Stacks, Tag~#1}}}
\providecommand{\stacksproject}{}\renewcommand{\stacksproject}[1]{{\cite[\href{https://stacks.math.columbia.edu/tag/#1}{Tag #1}]{stacks-project}}}
\providecommand{\statusnote}{}\renewcommand{\statusnote}[1]{\begin{quote}\small\raggedright\textbf{Status.} #1\end{quote}}
\providecommand{\tildeideal}{}\renewcommand{\tildeideal}[2]{\widetilde{(#1, #2)}}
\providecommand{\tildering}{}\renewcommand{\tildering}[3]{{#1}_{\widetilde{(#2, #3)}}}
\providecommand{\tr}{}\renewcommand{\tr}{\operatorname{tr}}
\providecommand{\wt}{}\renewcommand{\wt}{\operatorname{wt}}

\begin{document}
\pagestyle{plain}
\begin{center}
{\LARGE\bfseries Weil Conjectures\par}
\vspace{0.7em}
{\large Appendix A: Nagata compactification and Chow's lemma\par}
\vspace{0.5em}
{\small\color{bpmuted}Snapshot 2026-09-26, revision 8e4cb88e2a\par}
\end{center}
\vspace{0.5em}
{\small\tableofcontents}
\bigskip\hrule\bigskip

% object: chap-Nagata-compactification-and-Chows-lemma (blueprint/chap-Nagata-compactification-and-Chows-lemma/_section.md)
\section*{\texorpdfstring{A\quad Nagata compactification and Chow's lemma}{A Nagata compactification and Chow's lemma}}
\phantomsection\label{bp1}%
\addcontentsline{toc}{section}{\texorpdfstring{A\quad Nagata compactification and Chow's lemma}{A Nagata compactification and Chow's lemma}}
% text: prose-A-01
\phantomsection\label{bp2}%
The main text defines cohomology with compact support as in topology: to push forward along \(f:X\to Y\), embed \(X\) as an open subset of something proper over \(Y\), extend by zero and push forward, so that \(Rf_!=Rp_*j_!\) (Proposition \bpout{5.26}). This definition needs every morphism to have a compactification, and it needs any two compactifications to be dominated by a third, so that proper base change can show that \(Rf_!\) does not depend on the choice. The second point is elementary once we use compactifications in which \(X\) is scheme-theoretically dense: the closure of \(X\) in the fibre product of two of them dominates both (Proposition \hyperref[bp19]{A.7}). The first point is Nagata's theorem, Theorem \hyperref[bp7]{A.2}: every separated morphism of finite type \(X\to S\) with \(S\) quasi-compact and quasi-separated factors as a quasi-compact open immersion \(X\to\overline
X\) followed by a proper morphism \(\overline X\to S\).

Why is this hard? Locally a compactification is easy to find: if \(S\) and \(X\) are affine, we take the closure of \(X\) in projective space over \(S\). The difficulty is to glue. For example, cover \(X=\mathbf A^1_k\) by \(U_1=\mathbf A^1\setminus\{0\}\) and \(U_2=\mathbf A^1\setminus\{1\}\), compactify each by \(\mathbf P^1\), and glue the two copies of \(\mathbf
P^1\) along \(U_1\cap U_2\). The result is then \(\mathbf P^1\) with the three points \(0\), \(1\) and \(\infty\) doubled, which is not separated. For curves blow-ups are not needed, and they would not help either, since blowing up a point on a regular curve does nothing. In dimension at least two, however, the closures of \(U_1\setminus U_2\) and \(U_2\setminus U_1\) can meet in the boundary, and we pull them apart by blow-ups with centers in the boundary, which leave \(X\) unchanged (Lemma \hyperref[bp53]{A.20} and Proposition \hyperref[bp61]{A.23}). Induction over an affine cover then gives the theorem over a Noetherian base (Section \hyperref[bp63]{A.4}), and Noetherian approximation gives the general case, which we need because proper base change is proved over arbitrary qcqs bases.

Proper base change itself is computed on projective space, and Chow's lemma carries it to every proper morphism (Section \bpout{4.5}). We use it in the form of Theorem \hyperref[bp71]{A.25}: for a proper morphism \(X\to S\) there is a proper surjection \(\widetilde X\to X\) with \(\widetilde X\) projective over \(S\).

% object: sec-Compactifications (blueprint/chap-Nagata-compactification-and-Chows-lemma/sec-Compactifications/_section.md), content hidden
\subsection*{\texorpdfstring{A.1\quad Compactifications}{A.1 Compactifications}}
\phantomsection\label{bp3}%
\addcontentsline{toc}{subsection}{\texorpdfstring{A.1\quad Compactifications}{A.1 Compactifications}}
% text: prose-A-02
\phantomsection\label{bp4}%
To show that \(Rp_*j_!\) does not depend on the compactification, we must compare two of them, and in general neither maps to the other. The plane \(\mathbf A^2\) is open in \(\mathbf P^2\) and in \(\mathbf P^1\times\mathbf P^1\), but the identity of \(\mathbf A^2\) extends to a morphism in neither direction: the map \((x:y:z)\mapsto((x:z),(y:z))\) is undefined at \((1:0:0)\) and \((0:1:0)\), and its inverse is undefined at \((\infty,\infty)\). What exists is a third compactification mapping to both, here the closure of \(\mathbf A^2\) in \(\mathbf P^2\times\mathbf P^1\times\mathbf P^1\). So we compare two compactifications through a common refinement. A refinement \(h\) is required to satisfy \(h^{-1}(X)=X\). This is what lets us compute \(Rh_*j'_!\) by restriction to \(X\) in the proof of Proposition \bpout{5.26}.

Without density, refinements are not unique. For example, the point \(X=\operatorname{Spec}k\) is open in the compactification \(\operatorname{Spec}k\sqcup\mathbf P^1_k\), which is not dense. Every automorphism of \(\mathbf P^1_k\) gives a refinement of it to itself. Dense compactifications do not have this problem: morphisms between them are unique, and with their refinements they form a cofiltered category (Proposition \hyperref[bp19]{A.7}).

% object: definition-compactification-refinement (blueprint/chap-Nagata-compactification-and-Chows-lemma/sec-Compactifications/definition-compactification-refinement/_section.md)
\phantomsection\label{bp5}%
\begin{bpstatement}{Definition~A.1}{ (Compactification and refinement)}
Let \(f:X\to S\) be a morphism. A compactification is a factorization

\[
X\xrightarrow{j}\overline X\xrightarrow{p}S
\]

with \(j\) a quasi-compact open immersion and \(p\) proper. We say that it is \emph{dense} if \(j\) is scheme-theoretically dense, i.e., if the map \(\mathcal O_{\overline X}\to j_*\mathcal O_X\) is injective. A refinement of compactifications is an \(S\)-morphism \(h:\overline X'\to\overline X\) with \(hj'=j\) and \(h^{-1}(X)=X\), under the specified identifications with \(X\).
\end{bpstatement}

% text: prose-A-03
\phantomsection\label{bp6}%
For example, \(\mathbf A^1_k\subset\mathbf P^1_k\) is a dense compactification of \(\mathbf A^1_k\to\operatorname{Spec}k\), since a regular function on an open of \(\mathbf P^1_k\) that vanishes on \(\mathbf A^1_k\) is zero. The theorem says that such a factorization always exists.

% object: theorem-nagata-s-compactification-theorem (blueprint/chap-Nagata-compactification-and-Chows-lemma/sec-Compactifications/theorem-nagata-s-compactification-theorem/_section.md)
\phantomsection\label{bp7}%
\begin{bpstatementit}{Theorem~A.2}{ (Nagata's compactification theorem)}
Let \(S\) be quasi-compact and quasi-separated. Every separated morphism \(f:X\to S\) of finite type has a dense compactification.
\end{bpstatementit}

% text: prose-A-04
\phantomsection\label{bp8}%
This is [Stacks, \href{https://stacks.math.columbia.edu/tag/0F41}{0F41}], and we give the proof in Section \hyperref[bp63]{A.4}. We need it in this generality because proper base change is proved over arbitrary qcqs bases; for schemes of finite type over \(k\) the Noetherian case of the proof suffices.

Density is a condition on structure sheaves. For instance, a nilpotent section supported on the boundary is invisible on \(X\), but it prevents \(X\) from being schematically dense. The condition does two things for us. First, morphisms between dense compactifications are unique and common refinements exist (Proposition \hyperref[bp19]{A.7}). Second, passing to a dense compactification does not change the direct image with proper support (Lemma \hyperref[bp15]{A.5}).

% object: lemma-schematic-closure-uniqueness (blueprint/chap-Nagata-compactification-and-Chows-lemma/sec-Compactifications/lemma-schematic-closure-uniqueness/_section.md)
\phantomsection\label{bp9}%
\begin{bpstatementit}{Lemma~A.3}{ (Schematic closure)}
Let \(i:U\to T\) be a quasi-compact immersion, with \(T\) quasi-separated. Its scheme-theoretic image \(\overline U\subseteq T\) is a closed subscheme, \(U\to\overline U\) is an open immersion, and \(U\) is scheme-theoretically dense in \(\overline U\). The formation of \(\overline U\) commutes with flat base change \(T'\to T\), in particular with restriction to opens of \(T\).
\end{bpstatementit}

% proof: proof/lemma-schematic-closure-uniqueness
\begin{proof}[{Proof}]\phantomsection\label{bp10}%
The scheme-theoretic image is the smallest closed subscheme through which \(i\) factors. Its ideal is \(\ker(\mathcal O_T\to i_*\mathcal O_U)\), which is quasi-coherent because \(i\) is quasi-compact and quasi-separated. Its formation commutes with flat base change \(T'\to T\): pushforward along the quasi-compact quasi-separated \(i\) commutes with flat base change, and pullback along a flat morphism is exact, so the pullback of the ideal is the ideal of the base change [Stacks, \href{https://stacks.math.columbia.edu/tag/081I}{081I}]. Restriction to an open of \(T\) is a special case. On an open where the immersion is closed the image is therefore \(U\) itself, and this shows that \(U\) is open in its image. By construction, a section of \(\mathcal O_{\overline U}\) vanishing on \(U\) is zero, which is the definition of schematic density. For these statements see [Stacks, \href{https://stacks.math.columbia.edu/tag/01R8}{01R8}, \href{https://stacks.math.columbia.edu/tag/01RG}{01RG}].
\end{proof}

% text: prose-A-05
\phantomsection\label{bp11}%
Density is what makes extensions from \(U\) unique.

% object: lemma-uniqueness-maps-dense-open (blueprint/chap-Nagata-compactification-and-Chows-lemma/sec-Compactifications/lemma-uniqueness-maps-dense-open/_section.md)
\phantomsection\label{bp12}%
\begin{bpstatementit}{Lemma~A.4}{ (Uniqueness of extensions from a dense open)}
Let \(T\) be an \(S\)-scheme, let \(U\subseteq T\) be a scheme-theoretically dense open, and let \(Y\to S\) be separated. Two \(S\)-maps \(T\to Y\) that agree on \(U\) are equal.
\end{bpstatementit}

% proof: proof/lemma-uniqueness-maps-dense-open
\begin{proof}[{Proof}]\phantomsection\label{bp13}%
The equalizer of the two maps is a closed subscheme of \(T\), since it is the inverse image of the diagonal of \(Y/S\). Its ideal restricts to zero on \(U\), because the two maps agree there. As a section of \(\mathcal O_T\) that vanishes on a scheme-theoretically dense open is zero, the ideal is zero and the equalizer is all of \(T\).
\end{proof}

% text: prose-A-06
\phantomsection\label{bp14}%
By Lemma \hyperref[bp9]{A.3}, every compactification contains a dense one, the schematic closure of \(X\). For this replacement to be harmless we need that it does not change \(Rp_*j_!\).

% object: lemma-passing-schematic-closure (blueprint/chap-Nagata-compactification-and-Chows-lemma/sec-Compactifications/lemma-passing-schematic-closure/_section.md)
\phantomsection\label{bp15}%
\begin{bpstatementit}{Lemma~A.5}{ (Passing to the schematic closure)}
Let \(X\xrightarrow{j}\overline X\xrightarrow{p}S\) be any compactification with \(\overline X\) Noetherian, and let \(i:Z\hookrightarrow\overline X\) be the schematic closure of \(X\), with \(j_Z:X\hookrightarrow Z\) the induced open immersion. For a condensed coefficient ring \(\Lambda\), we use the pullbacks of its coefficient sheaf to these schemes as coefficients. For every \(K\in D(X_{\mathrm{pro}\acute{e}t},\Lambda_X)\) there are canonical identifications

\[
j_!K\simeq i_*j_{Z,!}K,\qquad
Rp_*j_!K\simeq R(p\circ i)_*j_{Z,!}K.
\]
\end{bpstatementit}

% proof: proof/lemma-passing-schematic-closure
\begin{proof}[{Proof}]\phantomsection\label{bp16}%
The first identification is a statement about supports, and the second follows from it by applying \(Rp_*\). The only point that needs care on the pro-étale site is the exactness of \(i_*\).

Since \(j_!K\) vanishes outside \(X\subseteq Z\), it should be pushed forward from \(Z\). To say this in the derived category we need \(i_*\) to be exact, so that \(Ri_*=i_*\). By Proposition \bpout{2.58}, the functors attached to open and closed immersions on the pro-étale site are exact when the complementary open is quasi-compact. Since \(\overline X\) is Noetherian, the complement of \(Z\) is quasi-compact, so all functors below are exact and may be applied degreewise to unbounded complexes. The inverse image of \(j_!K\) on \(Z\) is \(j_{Z,!}K\), so we have the natural map \(j_!K\to i_*i^*j_!K=i_*j_{Z,!}K\). It is the identity on the open \(X\), and both sides restrict to zero on the closed complement of \(X\). By the localization sequence for an open and its closed complement, this map is therefore an isomorphism. The second identification then follows by applying \(Rp_*\) and using \(Rp_*Ri_*\simeq R(p\circ i)_*\).
\end{proof}

% object: remark-non-noetherian-compactifications (blueprint/chap-Nagata-compactification-and-Chows-lemma/sec-Compactifications/remark-non-noetherian-compactifications/_section.md)
\phantomsection\label{bp17}%
\begin{bpstatement}{Remark~A.6}{ (Non-Noetherian compactifications)}
The same proof works for qcqs \(\overline X\) whenever \(\overline X\setminus Z\) is quasi-compact. For a non-Noetherian compactification this can fail, and then Proposition \bpout{2.58} does not apply to pushforward along the closed immersion on the pro-étale site. The schematic closure and dense refinements exist in the generality of Theorem \hyperref[bp7]{A.2}.
\end{bpstatement}

% text: prose-A-07
\phantomsection\label{bp18}%
For dense compactifications the comparison is formal. We obtain common refinements by Lemma \hyperref[bp9]{A.3}, and morphisms are unique by Lemma \hyperref[bp12]{A.4}.

% object: proposition-dense-compactifications-form-cofiltered-category (blueprint/chap-Nagata-compactification-and-Chows-lemma/sec-Compactifications/proposition-dense-compactifications-form-cofiltered-category/_section.md)
\phantomsection\label{bp19}%
\begin{bpstatementit}{Proposition~A.7}{ (Dense compactifications form a cofiltered category)}
Let \(S\) be quasi-compact and quasi-separated, and let \(f:X\to S\) be separated and of finite type. Any finite family of dense compactifications of \(f\) has a common dense refinement. Morphisms between dense compactifications that restrict to the identity on \(X\) are unique, proper, and automatically satisfy \(h^{-1}(X)=X\). The category of dense compactifications of \(f\) is therefore cofiltered.
\end{bpstatementit}

% proof: proof/proposition-dense-compactifications-form-cofiltered-category
\begin{proof}[{Proof}]\phantomsection\label{bp20}%
A common refinement of two compactifications is the closure of \(X\) in their fibre product, as in the example of \(\mathbf A^2\) above. The properties of morphisms then follow from density.

Given two compactifications, let \(Z\) be the schematic closure of the image of \(X\) in \(\overline X_1\times_S\overline X_2\) under the diagonal map. The map from \(X\) is a quasi-compact immersion, because over the open \(X\subseteq\overline X_1\) it is the graph of a morphism to \(\overline X_2\), and \(\overline X_2\) is separated over \(S\). By Lemma \hyperref[bp9]{A.3}, \(X\) is an open subscheme of \(Z\) and is scheme-theoretically dense in \(Z\). Since the product is proper over \(S\), the scheme \(Z\) and its projections are proper.

The projections are refinements, that is, \(X\) is its own inverse image in \(Z\). By Lemma \hyperref[bp9]{A.3} the formation of the schematic closure commutes with restriction to opens. Over \(X\subseteq\overline X_1\) we may therefore compute it on the open \(X\times_S\overline X_2\). There the immersion is the graph of \(X\to\overline X_2\), which is closed, so the closure is this graph. This means that the inverse image of \(X\) under the first projection is \(X\), on which the first projection is the identity. The same holds for the second projection. The assertion for finite families follows by induction.

Let \(h:\overline X'\to\overline X\) be a morphism of dense compactifications restricting to the identity on \(X\). To see that \(h\) is proper, note that it factors through its graph. This graph is closed in \(\overline X'\times_S\overline X\) because \(\overline X\) is separated over \(S\), and the projection of this product to \(\overline X\) is proper because \(\overline X'\) is proper over \(S\).

Next, the inclusion \(X\to h^{-1}(X)\) is a section of the separated map \(h^{-1}(X)\to X\). It is therefore closed, as well as open. Since \(X\) is schematically dense in \(\overline X'\), it is schematically dense in the open \(h^{-1}(X)\). A closed subscheme containing it must be all of \(h^{-1}(X)\), and we conclude that \(h^{-1}(X)=X\). Two parallel morphisms are equal by Lemma \hyperref[bp12]{A.4}. Finally, the category is nonempty by Nagata's theorem. Since any two objects have a common refinement and parallel morphisms are equal, it is cofiltered.
\end{proof}

% text: prose-A-08
\phantomsection\label{bp21}%
Proper base change is applied after a change of base \(S'\to S\), so we also need to know what happens to a compactification under base change, where density need not survive.

% object: remark-base-change-compactification (blueprint/chap-Nagata-compactification-and-Chows-lemma/sec-Compactifications/remark-base-change-compactification/_section.md)
\phantomsection\label{bp22}%
\begin{bpstatement}{Remark~A.8}{ (Base change of a compactification)}
For every \(S'\to S\), the base change of a compactification is an open immersion followed by a proper morphism, which suffices to define the direct image with compact support of \(X\times_SS'\to S'\). When the base-changed ambient scheme is Noetherian, as for schemes of finite type over \(k\), Lemma \hyperref[bp15]{A.5} shows that this direct image agrees with the one defined by the schematic closure.

Schematic closure commutes with flat base change (Lemma \hyperref[bp9]{A.3}), and so does schematic density, but density can be lost under an arbitrary base change. In that case the schematic closure of \(X\times_SS'\) in \(\overline X\times_SS'\) is a dense compactification, and a base-changed common refinement can be replaced by its closure in the same way. The comparisons in the main text are made on a further common refinement, where they are canonical by uniqueness.
\end{bpstatement}


% object: sec-Elimination-of-indeterminacy (blueprint/chap-Nagata-compactification-and-Chows-lemma/sec-Elimination-of-indeterminacy/_section.md), content hidden
\subsection*{\texorpdfstring{A.2\quad Elimination of indeterminacy}{A.2 Elimination of indeterminacy}}
\phantomsection\label{bp23}%
\addcontentsline{toc}{subsection}{\texorpdfstring{A.2\quad Elimination of indeterminacy}{A.2 Elimination of indeterminacy}}
% text: prose-A-09
\phantomsection\label{bp24}%
Every modification in the proof of Nagata's theorem is a blow-up whose center avoids a fixed open, so that this open is not changed. Here the open \(U\) is the scheme being compactified, or a piece of it, and only the boundary may be modified. For instance, the coordinate axes minus the origin are disjoint closed subsets of \(U=\mathbf A^2\setminus\{0\}\), but their closures in \(\mathbf A^2\) meet at the origin. The blow-up of the origin is an isomorphism over \(U\), and in it the strict transforms of the axes meet the exceptional line in the two distinct points of slope \(0\) and \(\infty\), so they become disjoint. Lemma \hyperref[bp53]{A.20} proves the same in general.

We use blow-ups rather than arbitrary proper modifications that are isomorphisms over \(U\), because a blow-up is given by an ideal. This makes them easy to combine. Two of them are dominated by the blow-up of the product ideal, and a composite is again a blow-up (Lemma \hyperref[bp33]{A.13}). The ideal is required to be of finite type, so that the blow-up is proper.

% object: definition-admissible-blow-up (blueprint/chap-Nagata-compactification-and-Chows-lemma/sec-Elimination-of-indeterminacy/definition-admissible-blow-up/_section.md)
\phantomsection\label{bp25}%
\begin{bpstatement}{Definition~A.9}{ (Admissible blow-up)}
Let \(T\) be quasi-compact and quasi-separated, and let \(U\subseteq T\) be a quasi-compact open. For a quasi-coherent ideal of finite type \(\mathcal I\subseteq\mathcal O_T\) with \(\mathcal I|_U=\mathcal O_U\), the blow-up

\[
b:\operatorname{Bl}_{\mathcal I}(T)=
\operatorname{Proj}_T\Bigl(\bigoplus_{m\geq0}\mathcal I^m\Bigr)
\longrightarrow T
\]

is called \(U\)-admissible. Its center is disjoint from \(U\).
\end{bpstatement}

% text: prose-A-10
\phantomsection\label{bp26}%
On an affine open \(\operatorname{Spec}(R)\), write \(I=(a_1,\ldots,a_r)\). The charts of the blow-up have rings

\[
R[I/a_i]=\bigl((\bigoplus_{m\geq0}I^mt^m)_{a_it}\bigr)_0.
\]

This is the degree-zero part of the localization at \(a_it\) of the Rees algebra, where \(t\) is a variable of degree one. It is the quotient of \(R[x_1,\ldots,x_r]/(a_ix_j-a_j)_j\) by its \(a_i\)-power torsion, that is, by the elements killed by some power of \(a_i\), and this torsion can be nonzero. In the chart \(R[I/a_i]\), the element \(a_i\) is a nonzerodivisor and generates the extended ideal \(IR[I/a_i]\). In other words, the center becomes an effective Cartier divisor. The blow-up is an isomorphism away from the center. Since \(I\) is finitely generated, the blow-up of an affine base is a closed subscheme of a projective space over it, and therefore \(b\) is proper. These constructions glue under restriction to opens.

We first use admissible blow-ups to restore density. A compactification may contain components or embedded components supported in the boundary. Blowing up the boundary removes them, because after the blow-up the boundary is cut out locally by a nonzerodivisor.

% object: lemma-making-open-schematically-dense (blueprint/chap-Nagata-compactification-and-Chows-lemma/sec-Elimination-of-indeterminacy/lemma-making-open-schematically-dense/_section.md)
\phantomsection\label{bp27}%
\begin{bpstatementit}{Lemma~A.10}{ (Making an open schematically dense)}
Let \(T\) be quasi-compact and quasi-separated, and let \(U\subseteq T\) be a quasi-compact open. There is a \(U\)-admissible blow-up \(T'\to T\) such that the inverse image of \(U\) is schematically dense in \(T'\) and its complement is the support of an effective Cartier divisor.
\end{bpstatementit}

% proof: proof/lemma-making-open-schematically-dense
\begin{proof}[{Proof}]\phantomsection\label{bp28}%
We blow up an ideal that cuts out the complement of \(U\). Its pullback is then locally generated by a nonzerodivisor \(a\), and an open whose complement is cut out by a nonzerodivisor is schematically dense, because a section killed by a power of \(a\) is zero.

Choose a quasi-coherent ideal \(\mathcal I\) of finite type with \(V(\mathcal I)=T\setminus U\). Such an ideal exists because \(U\) is a quasi-compact open in a qcqs scheme [Stacks, \href{https://stacks.math.columbia.edu/tag/01PH}{01PH}]. We blow up this ideal, and the extended ideal defines the exceptional Cartier divisor. On the affine charts its local equation \(a\) is a nonzerodivisor. Hence each chart ring injects into its localization at \(a\). The complement of the divisor is the inverse image of \(U\), and it is therefore schematically dense.
\end{proof}

% object: example-blow-up-embedded-point (blueprint/chap-Nagata-compactification-and-Chows-lemma/sec-Elimination-of-indeterminacy/example-blow-up-embedded-point/_section.md)
\phantomsection\label{bp29}%
\begin{bpstatement}{Example~A.11}{ (Blowing up removes an embedded point)}
Let \(T=\operatorname{Spec}R\) with \(R=k[x,y]/(xy,y^2)\), the line \(y=0\) with an embedded point at the origin, and let \(U=D(x)\). On \(U\) the element \(y\) is zero, since \(xy=0\) and \(x\) is invertible, so \(U\) is not schematically dense in \(T\). On the other hand, the ideal \(I=(x)\) is the unit ideal on \(U\), so its blow-up is \(U\)-admissible. Since the ideal is principal, the chart formula above gives \(R[I/x]=R/\{b:x^nb=0\text{ for some }n\}\). Now every element of \(R\) is uniquely \(f(x)+cy\) with \(c\in k\), and \(x^n(f+cy)=x^nf\), so the \(x\)-power torsion is \(ky\) and the blow-up is \(\operatorname{Spec}k[x]\). Thus the embedded point is gone, \(U\) is dense, and the boundary is the Cartier divisor \(x=0\), as in Lemma \hyperref[bp27]{A.10}.
\end{bpstatement}

% text: prose-A-11
\phantomsection\label{bp30}%
After a blow-up, we replace a scheme over \(T\) by the part of its base change that is not concentrated on the exceptional divisor.

% object: definition-strict-transform (blueprint/chap-Nagata-compactification-and-Chows-lemma/sec-Elimination-of-indeterminacy/definition-strict-transform/_section.md)
\phantomsection\label{bp31}%
\begin{bpstatement}{Definition~A.12}{ (Strict transform)}
Let \(b:T'\to T\) be the blow-up of \(T\) in a quasi-coherent ideal of finite type, with exceptional divisor \(E\subseteq T'\), and let \(W\to T\) be a morphism. The \emph{strict transform} of \(W\) is the closed subscheme of \(W\times_TT'\) whose ideal consists of the local sections of \(\mathcal O_{W\times_TT'}\) supported on the inverse image of \(E\).
\end{bpstatement}

% text: prose-A-12
\phantomsection\label{bp32}%
Over a chart of \(T'\) where \(E\) has equation \(a\) and \(W\times_TT'\) is affine with ring \(B\), the ring of the strict transform is

\[
B/\{b\in B:\ a^nb=0\text{ for some }n\geq0\}.
\]

If \(U\) is the complement of the center, this is exactly the schematic closure of \(W_U\) in \(W\times_TT'\), the object that occurs in the elimination of indeterminacy. By [Stacks, \href{https://stacks.math.columbia.edu/tag/080E}{080E}], the strict transform is the blow-up of \(W\) in the pulled-back ideal, so it is proper over \(W\).

In the proofs below we construct blow-ups locally and one after another. By the next lemma, finitely many such constructions combine into a single \(U\)-admissible blow-up of \(T\).

% object: lemma-products-composites-u-admissible-blow (blueprint/chap-Nagata-compactification-and-Chows-lemma/sec-Elimination-of-indeterminacy/lemma-products-composites-u-admissible-blow/_section.md)
\phantomsection\label{bp33}%
\begin{bpstatementit}{Lemma~A.13}{ (Products and composites of $U$-admissible blow-ups)}
Finitely many \(U\)-admissible blow-ups of \(T\) admit a common \(U\)-admissible blow-up. A composite of two \(U\)-admissible blow-ups is again such a blow-up. A \(U\)-admissible blow-up on a quasi-compact open \(V\subseteq T\) containing \(U\) extends to a \(U\)-admissible blow-up of \(T\).
\end{bpstatementit}

% proof: proof/lemma-products-composites-u-admissible-blow
\begin{proof}[{Proof}]\phantomsection\label{bp34}%
All three assertions are statements about ideals, because products of ideals dominate blow-ups, composites are computed by products of ideals, and extension from an open amounts to extending a quasi-coherent ideal of finite type. Let \(\mathcal I,\mathcal J\) be ideals, both equal to \(\mathcal O\) on \(U\), so that their product is too. Then the blow-up of their product factors through both blow-ups. Indeed, on it the product \(\mathcal I\mathcal J\) is invertible, hence so is each factor, and the universal property of each blow-up then gives the factorization. The first assertion follows by induction. For the second, blowing up first \(\mathcal I\) and then the pullback of \(\mathcal J\) yields the blow-up of \(\mathcal I\mathcal J\) [Stacks, \href{https://stacks.math.columbia.edu/tag/080A}{080A}], which proves the composition assertion when the second center is pulled back from \(T\). In general, however, the second center lives only on the first blow-up, and the composite is a blow-up of \(T\) by the calculation with ideals of finite type in [Stacks, \href{https://stacks.math.columbia.edu/tag/080B}{080B}]. For the extension assertion we extend a quasi-coherent ideal of finite type from a quasi-compact open, keeping the center outside \(U\); such an extension exists because on a qcqs scheme every quasi-coherent sheaf is the union of its quasi-coherent subsheaves of finite type, and the precise result on admissible extensions is [Stacks, \href{https://stacks.math.columbia.edu/tag/080M}{080M}].
\end{proof}

% text: prose-A-13
\phantomsection\label{bp35}%
Blow-ups commute with flat base change [Stacks, \href{https://stacks.math.columbia.edu/tag/0805}{0805}], because flat pullback preserves the inclusions \(\mathcal I^m\subseteq\mathcal O_T\). Under a base change that is not flat, the pullbacks of these inclusions can have kernels. In that case the blow-up of the pulled-back ideal is the strict transform of the base change, which can be strictly smaller than the base change itself.

The gluing argument has to extend a morphism defined on a dense open \(U\subseteq X\) as far as possible. When the target \(Y\) is not proper, the extension can stop before reaching all of \(X\). We need a way to say that a partial map is already extended as far as it can be.

Take \(X=Y=\mathbf A^1_k\) and \(U=\mathbf A^1\setminus\{0\}\). The inclusion \(t\mapsto t\) extends over \(0\). Its graph, the diagonal minus the origin, is not closed in \(X\times Y\). The map \(t\mapsto1/t\) does not extend, even after a blow-up, since blowing up a point on a regular curve does nothing. Its graph \(V(ty-1)\) is closed. In general, a closed graph means that no limit point lies over \(X\setminus U\). A map with closed graph is therefore not the restriction of a map on a larger open. Both the hypothesis and the conclusion of Theorem \hyperref[bp41]{A.16} are stated in this form. For proper \(Y\) the condition forces the domain to be all of \(X\), which gives Corollary \hyperref[bp43]{A.17}.

% object: definition-quasi-domination (blueprint/chap-Nagata-compactification-and-Chows-lemma/sec-Elimination-of-indeterminacy/definition-quasi-domination/_section.md)
\phantomsection\label{bp36}%
\begin{bpstatement}{Definition~A.14}{ (Quasi-domination)}
Let \(X\) and \(Y\) be \(S\)-schemes, with \(Y\to S\) separated. A \emph{quasi-domination of \(X\) over \(Y\)} consists of a dense open \(U\subseteq X\) and an \(S\)-map \(f:U\to Y\) whose graph is closed in \(X\times_SY\).
\end{bpstatement}

% text: prose-A-14
\phantomsection\label{bp37}%
Since \(Y\to S\) is separated, the graph is always closed in \(U\times_SY\); the condition is about its behaviour over \(X\setminus U\). Every map \(X\to Y\) is a quasi-domination with domain \(X\). Conversely, if \(Y\to S\) is proper, every quasi-domination has domain \(X\), because the projection of the closed graph to \(X\) is closed with image the dense open \(U\).

The extension theorem below is proved first for targets that embed in affine space over the base. We reduce a general target to this form by Zariski's Main Theorem.

% object: proposition-relative-form-zariski-s-main (blueprint/chap-Nagata-compactification-and-Chows-lemma/sec-Elimination-of-indeterminacy/proposition-relative-form-zariski-s-main/_section.md)
\phantomsection\label{bp38}%
\begin{bpstatementit}{Proposition~A.15}{ (The relative form of Zariski's Main Theorem)}
A separated, quasi-compact, locally quasi-finite morphism \(Z\to T\) locally of finite type factors as an open immersion \(Z\hookrightarrow T_0\) followed by an affine morphism \(T_0\to T\). Consequently \(Z\to T\) is quasi-affine.
\end{bpstatementit}

% proof: proof/proposition-relative-form-zariski-s-main
\begin{proof}[{Proof}]\phantomsection\label{bp39}%
The scheme \(T_0\) is the normalization of \(T\) in \(Z\), that is, the relative spectrum of the integral closure of \(\mathcal O_T\) in the direct image of \(\mathcal O_Z\); it is affine over \(T\) by construction. By Zariski's Main Theorem in the form [Stacks, \href{https://stacks.math.columbia.edu/tag/02LR}{02LR}], locally on \(T\) the canonical map from \(Z\) to the normalization \(T_0\) of \(T\) in \(Z\) is an open immersion, and \(T_0\to T\) is affine. These factorizations agree on overlaps and glue to the stated factorization. Over an affine open of \(T\) the scheme \(T_0\) is affine, so \(Z\) is a quasi-compact open of an affine scheme there, which means that \(Z\to T\) is quasi-affine.
\end{proof}

% text: prose-A-15
\phantomsection\label{bp40}%
Quasi-affine targets are the ones where a partial section can be extended by hand, by clearing denominators of finitely many coordinates. Using the proposition, we reduce the general extension theorem to this case, once we know that the relevant closure is quasi-finite over the base.

% object: theorem-extension-partial-section (blueprint/chap-Nagata-compactification-and-Chows-lemma/sec-Elimination-of-indeterminacy/theorem-extension-partial-section/_section.md)
\phantomsection\label{bp41}%
\begin{bpstatementit}{Theorem~A.16}{ (Extension of a partial section)}
Let \(T\) be qcqs, let \(U\subseteq V\subseteq T\) be quasi-compact opens with \(U\) dense in \(T\), and let \(p:Y\to T\) be separated and of finite type. Suppose that \(f:U\to Y\) is a section of \(p\) over \(U\) and that its image is closed in \(p^{-1}(V)\). There are a \(V\)-admissible blow-up \(\pi:T'\to T\), an open \(W\subseteq T'\) containing \(\pi^{-1}(U)\), and a map \(g:W\to Y\) such that

\[
p\circ g=\pi|_W,\qquad
g|_{\pi^{-1}(U)}=f\circ\pi|_{\pi^{-1}(U)},
\]

and the graph of \(g\) is closed in \(Y\times_TT'\).
\end{bpstatementit}

% text: prose-A-16
\phantomsection\label{bp42}%
We quote this theorem from [\hyperref[bp73]{Con07}, Theorem 2.4], with the auxiliary base equal to \(T\), and explain below the structure of its proof. Note that \(Y\) is qcqs, since \(p\) is separated and of finite type and \(T\) is qcqs. The hypothesis on \(f\) says that \(f\) is a quasi-domination of \(V\) over \(Y\). We apply the theorem only when \(V=U\), the case of [\hyperref[bp73]{Con07}, Remark 2.5]. Then the hypothesis holds automatically, because a section of a separated morphism is a closed immersion.

The proof first treats a quasi-affine \(Y\to T\). If \(Y\to T\) is quasi-affine and \(T\) is affine, then \(Y\) is a quasi-affine scheme of finite type over \(T\), and it embeds in some \(\mathbf A^n_T\) [Stacks, \href{https://stacks.math.columbia.edu/tag/04II}{04II}]. The coordinates of \(f\) are then finitely many fractions, and a blow-up clearing their denominators extends \(f\). One first makes \(U\) schematically dense, so that the extended coordinates still satisfy the equations of \(Y\). Using a finite affine cover of \(T\), we then obtain the quasi-affine case over a qcqs base. In fact the proof needs this case for opens \(U\subseteq V\) with \(U\) not necessarily dense, because the reduction below applies it to the opens \(f^{-1}(Y_i)\subseteq V\), which need not be dense. For this, a further construction in Conrad's proof keeps the centers outside \(V\).

The general case reduces to the quasi-affine one, and this reduction is where Zariski's Main Theorem enters. The plan is to extend the section on each affine open of \(Y\) separately. The closure of the section is then quasi-finite over the base, hence quasi-affine, and the quasi-affine case applies once more.

In detail, choose a finite affine open cover \(Y=\bigcup_iY_i\). Since \(T\) is quasi-separated, each \(Y_i\to T\) is quasi-compact, and as \(Y_i\) is affine it is quasi-affine. Hence the quasi-affine case applies to each \(f|_{f^{-1}(Y_i)}\), and we take a common \(V\)-admissible blow-up \(T_1\to T\). Let \(Z\) be the schematic closure of the image of \(f\) in \(Y\times_TT_1\). It meets each \(Y_i\times_TT_1\) inside the closed graph of an extension, which has at most one point over each point of \(T_1\). Since there are finitely many \(Y_i\), the fibres of \(Z\to T_1\) are finite. As \(Z\to T_1\) is separated and of finite type, it is quasi-finite, and it is quasi-affine by Proposition \hyperref[bp38]{A.15}. The quasi-affine case then extends the section into \(Z\), and the composite of the two blow-ups is \(V\)-admissible by Lemma \hyperref[bp33]{A.13}.

% object: corollary-extension-proper-target (blueprint/chap-Nagata-compactification-and-Chows-lemma/sec-Elimination-of-indeterminacy/corollary-extension-proper-target/_section.md)
\phantomsection\label{bp43}%
\begin{bpstatementit}{Corollary~A.17}{ (Extension to a proper target)}
Let \(X\) be qcqs over \(S\), let \(U\subseteq X\) be a dense quasi-compact open and let \(Y\to S\) be proper. Every \(S\)-map \(f:U\to Y\) extends to an \(S\)-map \(X'\to Y\) for some \(U\)-admissible blow-up \(X'\to X\).
\end{bpstatementit}

% proof: proof/corollary-extension-proper-target
\begin{proof}[{Proof}]\phantomsection\label{bp44}%
A map \(U\to Y\) is the same as a section over \(U\) of the projection \(Y\times_SX\to X\). The theorem extends this section to an open of a blow-up, and properness of \(Y\) forces this open to be everything.

Apply Theorem \hyperref[bp41]{A.16}, with \(T=X\) and \(V=U\), to the section \((f,\mathrm{id}):U\to Y\times_SX\) of the proper projection to \(X\). This gives an extension on an open \(W\subseteq X'\) whose graph is closed in \(Y\times_SX'\). The open \(W\) is dense, because it contains the inverse image of \(U\), which is dense in \(X'\). Since \(Y\to S\) is proper, the observation after Definition \hyperref[bp36]{A.14} shows that \(W=X'\).
\end{proof}

% text: prose-A-17
\phantomsection\label{bp45}%
When the target is not proper, the extension is defined only on an open of the blow-up, and we can still say what it looks like. Suppose the target is a separated map \(W\to T\) of finite type that is an isomorphism over \(U\), and the section is its inverse over \(U\). Then the extension identifies a blow-up of \(W\) with an open of a blow-up of \(T\).

% object: theorem-elimination-indeterminacy (blueprint/chap-Nagata-compactification-and-Chows-lemma/sec-Elimination-of-indeterminacy/theorem-elimination-indeterminacy/_section.md)
\phantomsection\label{bp46}%
\begin{bpstatementit}{Theorem~A.18}{ (Elimination of indeterminacy)}
Let \(T\) be qcqs, let \(U\subseteq T\) be a quasi-compact open, and let \(W\to T\) be separated and of finite type, with \(W_U\to U\) an isomorphism. There is a \(U\)-admissible blow-up \(T'\to T\) such that the inverse image of \(U\) is schematically dense in \(T'\) and the strict transform \(W'\to T'\) of \(W\to T\) is an open immersion.
\end{bpstatementit}

% proof: proof/theorem-elimination-indeterminacy
\begin{proof}[{Proof}]\phantomsection\label{bp47}%
This is [Stacks, \href{https://stacks.math.columbia.edu/tag/081S}{081S}]. The idea is to view the inverse of \(W_U\to U\) as a partial section of \(W\to T\) and to extend it by Theorem \hyperref[bp41]{A.16}; the strict transform is then the graph of the extension.

Choose a finite type ideal \(\mathcal I_0\) cutting out \(T\setminus U\). Its blow-up \(T_1\to T\) makes \(U\) schematically dense by Lemma \hyperref[bp27]{A.10}. Over \(T_1\), the base change of \(W\to T\) to \(T_1\) is separated and of finite type, and the inverse of \(W_U\to U\) is a section of it over \(U\). We may therefore apply Theorem \hyperref[bp41]{A.16} over \(T_1\), with \(V=U\), which gives a further blow-up and an extension of this section, with closed graph, over an open containing the inverse image of \(U\). By Lemma \hyperref[bp33]{A.13} the composite is a \(U\)-admissible blow-up \(\operatorname{Bl}_{\mathcal I}(T)\). Let \(b:T'\to T\) be the blow-up of \(\mathcal I\mathcal I_0\), which dominates both. Pulling back to \(T'\) gives an open \(W_0\subseteq T'\) containing \(b^{-1}(U)\) and a map \(W_0\to W\) over \(T\) whose graph \(\Gamma\) is closed in \(W\times_TT'\). As the zero set of \(\mathcal I\mathcal I_0\) is exactly \(T\setminus U\), \(b^{-1}(U)\) is schematically dense in \(T'\), hence in \(W_0\).

We included the factor \(\mathcal I_0\) so that the center of \(b\) is exactly \(T\setminus U\). Then the strict transform \(W'\) is the schematic closure of \(W_U\) in \(W\times_TT'\). Over \(U\), the subscheme \(W_U\) is the graph of the original section, so it lies in \(\Gamma\) and \(W'\subseteq\Gamma\). Conversely, \(\Gamma\simeq W_0\) contains the inverse image of \(U\) as a schematically dense open. Hence the schematic closure of \(W_U\) in \(\Gamma\) is \(\Gamma\). We conclude that \(W'=\Gamma\simeq W_0\), and \(W'\to T'\) is an open immersion.
\end{proof}

% text: prose-A-18
\phantomsection\label{bp48}%
The blow-up is not uniquely determined. However, on any such modification with \(U\) schematically dense, the graph of the extension is forced to be the closure of the original section. By Lemma \hyperref[bp12]{A.4}, the identifications obtained on common admissible refinements therefore agree.

In the gluing argument below, a compactification of an open \(V\) of \(M\) is related to \(M\) itself by a proper map defined on a larger open of the compactification.

% object: lemma-extending-graph-from-compactification (blueprint/chap-Nagata-compactification-and-Chows-lemma/sec-Elimination-of-indeterminacy/lemma-extending-graph-from-compactification/_section.md)
\phantomsection\label{bp49}%
\begin{bpstatementit}{Lemma~A.19}{ (Extending a graph from a compactification)}
Let \(S\) be qcqs, let \(M\to S\) be separated of finite type, and let \(V\subseteq M\) be a quasi-compact open with a compactification \(V\subseteq Y\) over \(S\). There are a \(V\)-admissible blow-up \(Y'\to Y\), an open \(V\subseteq V'\subseteq Y'\) and a proper map \(\psi:V'\to M\) with \(\psi^{-1}(V)=V\) and \(\psi|_V=\mathrm{id}_V\).
\end{bpstatementit}

% proof: proof/lemma-extending-graph-from-compactification
\begin{proof}[{Proof}]\phantomsection\label{bp50}%
The map \(\psi\) is the projection to \(M\) of the closure of the graph of \(V\to M\). By elimination of indeterminacy this closure becomes an open of a blow-up of \(Y\), and \(\psi\) is proper because the graph is closed.

Let \(G\) be the schematic closure of the graph of \(V\to M\) in \(Y\times_SM\). Its projection to \(Y\) is separated of finite type, and it is an isomorphism over \(V\). By Theorem \hyperref[bp46]{A.18}, after a \(V\)-admissible blow-up \(Y'\to Y\) its strict transform \(G'\) is an open \(V'\subseteq Y'\). This strict transform is the closure of the graph of \(V\to M\) in \(Y'\times_SM\), since both are the closure of the same graph outside the exceptional divisor [Stacks, \href{https://stacks.math.columbia.edu/tag/0F3Y}{0F3Y}]. So \(G'\) is closed in \(Y'\times_SM\), and since the projection of \(Y'\times_SM\) to \(M\) is proper, the map \(V'\simeq G'\to M\) is proper. It remains to identify the inverse image of \(V\). For this, note that over \(V\subseteq M\) the graph of the inclusion \(V\to Y'\) is already closed in \(Y'\times_SV\), because \(Y'\) is separated over \(S\). The inverse image of \(V\) under \(\psi\) is therefore this graph, that is, \(V\) itself. See also [Stacks, \href{https://stacks.math.columbia.edu/tag/0F3Z}{0F3Z}].
\end{proof}


% object: sec-Separation-by-blow-ups (blueprint/chap-Nagata-compactification-and-Chows-lemma/sec-Separation-by-blow-ups/_section.md), content hidden
\subsection*{\texorpdfstring{A.3\quad Separation by blow-ups}{A.3 Separation by blow-ups}}
\phantomsection\label{bp51}%
\addcontentsline{toc}{subsection}{\texorpdfstring{A.3\quad Separation by blow-ups}{A.3 Separation by blow-ups}}
% text: prose-A-19
\phantomsection\label{bp52}%
The inductive step of Nagata's theorem is Proposition \hyperref[bp61]{A.23}. Suppose that \(M=M_1\cup M_2\) with \(M_1\cap M_2\) dense, and that each \(M_i\) has a compactification \(X_i\); we want a compactification of \(M\). The natural candidate glues \(M\) to each \(X_i\) along \(M_i\), after we remove from \(X_i\) the closure of the part lying over \(M\setminus M_i\). For the result to be separated, however, and for every valuation to extend into one of the two pieces, the closures in \(X_i\) of the two disjoint closed subsets over \(M\setminus M_1\) and \(M\setminus M_2\) must be disjoint. They need not be, and this is where Lemma \hyperref[bp53]{A.20} enters: it makes them disjoint by an admissible blow-up, as with the coordinate axes in Section \hyperref[bp23]{A.2}.

We need two other lemmas. Lemma \hyperref[bp55]{A.21} keeps an isomorphism near such a closure, and Lemma \hyperref[bp58]{A.22} glues two modifications into one separated scheme. Both use elimination of indeterminacy from the previous section. Recall that it turns a map that is an isomorphism over a quasi-compact open \(U\) into an open immersion after a \(U\)-admissible blow-up.

% object: lemma-separation-closures (blueprint/chap-Nagata-compactification-and-Chows-lemma/sec-Separation-by-blow-ups/lemma-separation-closures/_section.md)
\phantomsection\label{bp53}%
\begin{bpstatementit}{Lemma~A.20}{ (Separation of closures)}
Let \(T\) be qcqs and let \(V\subseteq T\) be a quasi-compact open. If \(F_1,F_2\subseteq V\) are disjoint constructible closed subsets, there is a \(V\)-admissible blow-up of \(T\) on which their closures are disjoint.
\end{bpstatementit}

% text: prose-A-20
\phantomsection\label{bp54}%
We quote this from [Stacks, \href{https://stacks.math.columbia.edu/tag/0F3V}{0F3V}] and sketch the argument, which works with closed subschemes of \(T\) rather than closed subsets of \(V\). Let \(\mathcal I_1,\mathcal I_2\) be ideals of finite type on \(T\) cutting out closed subschemes \(Z_1,Z_2\) that are disjoint over \(V\) and restrict to \(F_1,F_2\) there. Producing such ideals on all of \(T\) is where constructibility and finite presentation are needed. Since \(\mathcal I_1+\mathcal I_2\) is the unit ideal over \(V\), its blow-up is \(V\)-admissible. Moreover, on each affine chart the extended ideal is generated by a single element of \(\mathcal I_1\) or of \(\mathcal I_2\). Now, on a chart where it is generated by \(f\in\mathcal I_1\), the strict transform of \(Z_1\) is obtained from its inverse image by dividing out the \(f\)-power torsion. Since \(f\) vanishes on \(Z_1\), every element of the ring of the inverse image is killed by \(f\). Hence the strict transform of \(Z_1\) does not meet that chart. The same holds for \(Z_2\) on the other charts, and therefore the strict transforms of \(Z_1\) and \(Z_2\) are disjoint [Stacks, \href{https://stacks.math.columbia.edu/tag/080P}{080P}]. The constructibility needed to produce such ideals causes no trouble for us, because in the Noetherian applications below every closed subset is constructible and its reduced ideal is of finite type.

For the gluing we also need to know that a proper map that is an isomorphism near a closed subset of \(V\) remains one near its closure in \(T\), after a blow-up.

% object: lemma-extending-isomorphism-over-closure-closed (blueprint/chap-Nagata-compactification-and-Chows-lemma/sec-Separation-by-blow-ups/lemma-extending-isomorphism-over-closure-closed/_section.md)
\phantomsection\label{bp55}%
\begin{bpstatementit}{Lemma~A.21}{ (Extending an isomorphism over the closure of a closed subset)}
Let \(q:P\to T\) be a proper morphism of qcqs schemes, let \(V\subseteq T\) be a quasi-compact open, and let \(F\subseteq V\) be closed. Suppose that \(q\) is an isomorphism over a neighbourhood of \(F\) in \(V\). After a \(V\)-admissible blow-up, the strict transform of \(q\) is an isomorphism over a neighbourhood of the closure of \(F\).
\end{bpstatementit}

% proof: proof/lemma-extending-isomorphism-over-closure-closed
\begin{proof}[{Proof}]\phantomsection\label{bp56}%
We separate the closure of \(F\) from the closure of the locus where \(q\) fails to be an isomorphism, and then apply elimination of indeterminacy on the complement of the second closure.

Inside \(V\), the set \(F\) is disjoint from the closed complement of the largest open over which \(q\) is an isomorphism. Since \(V\) is spectral, both closed subsets can be enlarged to disjoint constructible closed subsets. To these we apply Lemma \hyperref[bp53]{A.20} so that their closures in \(T\) become disjoint. Then, on the complement \(T_0\) of the second closure, the restriction of \(q\) is an isomorphism over an open \(V_0\) containing \(F\). After making this open schematically dense, we apply Theorem \hyperref[bp46]{A.18} on \(T_0\). The resulting strict transform is both open and proper over the modification, and hence it is an isomorphism there. By the extension part of Lemma \hyperref[bp33]{A.13}, this blow-up of \(T_0\) extends to a \(V\)-admissible blow-up of \(T\), which is an isomorphism over an open containing the closure of \(F\). We have only sketched the argument; the precise version, including the choice of constructible enlargements, is [Stacks, \href{https://stacks.math.columbia.edu/tag/0F3W}{0F3W}].
\end{proof}

% text: prose-A-21
\phantomsection\label{bp57}%
The last ingredient replaces two separated schemes containing a common open by modifications that glue to a separated scheme.

% object: lemma-gluing-after-admissible-blow-ups (blueprint/chap-Nagata-compactification-and-Chows-lemma/sec-Separation-by-blow-ups/lemma-gluing-after-admissible-blow-ups/_section.md)
\phantomsection\label{bp58}%
\begin{bpstatementit}{Lemma~A.22}{ (Gluing after admissible blow-ups)}
Let \(S\) be qcqs, and let \(U\subseteq T_i\) (\(i=1,2\)) be open immersions of separated \(S\)-schemes of finite type, with \(U\) quasi-compact. There are \(U\)-admissible blow-ups \(T_i'\to T_i\) and open immersions \(T_i'\to T\) agreeing on \(U\), with \(T\) separated and of finite type over \(S\). If \(U\) is dense in both \(T_i\), then it is dense in \(T\).
\end{bpstatementit}

% proof: proof/lemma-gluing-after-admissible-blow-ups
\begin{proof}[{Proof}]\phantomsection\label{bp59}%
We glue the \(T_i\) along the closure of \(U\) in \(T_1\times_ST_2\), after blow-ups that make this closure open in both. The glued scheme is separated because the closure is closed in the product.

Let \(G\) be the schematic closure of \(U\) in \(T_1\times_ST_2\). Then Theorem \hyperref[bp46]{A.18}, applied to the two projections, makes the strict transforms of \(G\) opens in suitable blow-ups of the two sides. To obtain a single modification, we blow up \(G\) in the product of the two pulled-back centers, which gives a common refinement. By Lemma \hyperref[bp33]{A.13}, these further centers extend from the opens to each side, and the resulting scheme \(G'\) is an open in both \(T_i'\).

The map \(G'\to T_1'\times_ST_2'\) is a closed immersion. It is an immersion, and it is proper because \(G'\to G\) is proper, \(G\subseteq T_1\times_ST_2\) is closed, and the modified product is separated over the original product. We glue the two \(T_i'\) along \(G'\), and the glued scheme \(T\) is covered by \(T_1'\) and \(T_2'\). To see that \(T\) is separated, note that the diagonal of \(T\) restricted to \(T_i'\times_ST_j'\) is the diagonal of \(T_i'\) for \(i=j\) and the graph \(G'\) for \(i\neq j\). Both are closed, so \(T\) is separated. It is of finite type because it has a finite cover by opens of finite type.

Suppose that \(U\) is dense in both \(T_i\). The complement of the exceptional divisor of a blow-up is schematically dense, as in the proof of Lemma \hyperref[bp27]{A.10}, and it contains the inverse image of \(U\) as a dense open. Hence \(U\) is dense in each \(T_i'\), and so in \(T\). This construction is [Stacks, \href{https://stacks.math.columbia.edu/tag/0F3X}{0F3X}].
\end{proof}

% text: prose-A-22
\phantomsection\label{bp60}%
The gluing lemma produces a separated scheme but says nothing about properness. In the inductive step we prove properness by the valuative criterion, using the separation of closures.

% object: proposition-compactifying-union-two-opens (blueprint/chap-Nagata-compactification-and-Chows-lemma/sec-Separation-by-blow-ups/proposition-compactifying-union-two-opens/_section.md)
\phantomsection\label{bp61}%
\begin{bpstatementit}{Proposition~A.23}{ (Compactifying a union of two opens)}
Let \(S\) be Noetherian and let \(M\to S\) be separated of finite type. Suppose that \(M=M_1\cup M_2\), that each \(M_i\) has a compactification over \(S\), and that \(D=M_1\cap M_2\) is dense in \(M\). Then \(M\) has a compactification over \(S\).
\end{bpstatementit}

% proof: proof/proposition-compactifying-union-two-opens
\begin{proof}[{Proof}]\phantomsection\label{bp62}%
The idea is to glue \(M\) to each compactification \(X_i\) of \(M_i\) after removing from \(X_i\) the closure of the part lying over \(M\setminus M_i\). This gives separated schemes \(W_i\) containing \(M\), and a common modification combines them into one separated \(W\). It remains to check that \(W\) is proper, and for this we use the valuative criterion. Separating the closures of the boundaries, we obtain (A.1), and from it we deduce that every valuation of \(D\) extends into one of the \(W_i\).

\emph{Separating the boundaries.} Since the subsets to be removed from \(X_i\) are defined by points of \(M\) outside \(M_i\), we must relate each \(X_i\) to \(M\) itself. Choose dense compactifications \(M_i\subseteq X_i\). By Lemma \hyperref[bp49]{A.19}, after admissible blow-ups there are opens \(M_i\subseteq V_i\subseteq X_i\) and proper maps \(\psi_i:V_i\to M\) that are isomorphisms above \(M_i\). Put, for \(\{i,j\}=\{1,2\}\),

\[
Z_i=M\setminus M_j,\qquad
Z_{i,r}=\psi_i^{-1}(Z_r),\qquad
B_{i,r}=\overline{Z_{i,r}}\subseteq X_i.
\]

The subsets \(Z_{i,1},Z_{i,2}\) are disjoint and closed in \(V_i\). By Lemma \hyperref[bp53]{A.20}, after a further \(V_i\)-admissible blow-up we may assume that \(B_{i,1}\cap B_{i,2}=\varnothing\), and subsequent \(V_i\)-admissible blow-ups preserve this disjointness.

Let \(X_{12}\) be the schematic closure of \(V_1\times_MV_2\) in \(X_1\times_SX_2\), with proper projections \(p_i\). Above \(V_1\), the first projection is the base change of \(\psi_2\), so it is an isomorphism near \(Z_{1,2}\). By Lemma \hyperref[bp55]{A.21} it becomes an isomorphism near \(B_{1,2}\) after a blow-up. We replace \(X_{12}\) by its strict transform, or equivalently by the closure of the new graph. Since \(p_1\) is now an isomorphism over a neighbourhood of \(B_{1,2}\) and \(Z_{1,2}\) is dense in \(B_{1,2}\), the inverse image of \(Z_{1,2}\) is dense in \(p_1^{-1}(B_{1,2})\). The second projection maps that inverse image into \(Z_{2,2}\). By density it maps all of \(p_1^{-1}(B_{1,2})\) into \(B_{2,2}\), which is disjoint from \(B_{2,1}\). This gives the disjointness used in the valuative argument below:

\[
X_{12}\cap(B_{1,2}\times_SB_{2,1})=\varnothing.
\tag{A.1}
\]

\emph{The separated schemes \(W_i\).} We glue

\[
W_i=M\coprod_{M_i}(X_i\setminus B_{i,j}).
\]

We claim that these schemes are separated over \(S\). The scheme \(W_i\) is covered by \(M\) and \(X_i\setminus B_{i,j}\), both separated, and they are glued along \(M_i\). It therefore suffices to show that the graph of the inclusion \(M_i\to X_i\setminus B_{i,j}\), which is the mixed part of the diagonal, is closed in \(M\times_S(X_i\setminus B_{i,j})\).

Suppose that a point of the graph of \(M_i\) specializes in \(M\times_S(X_i\setminus B_{i,j})\). Since \(\psi_i\) is proper, its graph is closed in \(M\times_SX_i\). Hence the specialized point lies in the graph of \(\psi_i\), over a point of \(V_i\). If its image in \(M\) lay in \(M\setminus M_i=Z_j\), the lifted point would lie in \(Z_{i,j}\subseteq B_{i,j}\), contradicting the second coordinate. So the specialization remains in the graph. The graph is an immersion with constructible image, and in the Noetherian situation a constructible subset that is closed under specialization is closed. Therefore the graph is a closed immersion, which proves separatedness.

\emph{Valuations of \(D\).} For properness we use the valuative criterion, and we show first that every valuation of \(D\) extends into one of the \(W_i\). Let \(R\) be a valuation ring with fraction field \(K\), and let \(\operatorname{Spec}(K)\to D\) be an \(S\)-map over a given \(\operatorname{Spec}(R)\to S\). It extends to both \(X_i\), since they are proper, and the pair of extensions factors through \(X_{12}\), because the closure of the graph contains its generic point and \(\operatorname{Spec}(K)\) is schematically dense in \(\operatorname{Spec}(R)\). By (A.1), the closed point cannot map into both of the removed closed subsets \(B_{1,2}\) and \(B_{2,1}\). We may therefore choose \(i\) such that the closed point does not map into \(B_{i,j}\). Since an open containing the closed point of a local scheme contains the whole scheme, the extension factors through \(X_i\setminus B_{i,j}\subseteq W_i\).

\emph{Gluing and properness.} Lemma \hyperref[bp58]{A.22}, applied to \(M\subseteq W_i\), gives a separated \(W\) of finite type, covered by admissible modifications \(W_i'\). Since the modification is proper, the lift of the valuation into \(W_i\) lifts further into \(W_i'\), hence into \(W\). Such lifts are unique because \(W\) is separated. The valuative criterion [Stacks, \href{https://stacks.math.columbia.edu/tag/0894}{0894}], in the form that only requires lifts of maps from \(\operatorname{Spec}(K)\) into the dense open \(D\), then shows that \(W\to S\) is universally closed. As \(W\) is also of finite type and separated, it is proper over \(S\). Its open \(M\) yields the desired compactification, which becomes dense after we take the schematic closure of \(M\). This is [Stacks, \href{https://stacks.math.columbia.edu/tag/0F40}{0F40}].
\end{proof}


% object: sec-Proof-of-Nagatas-theorem (blueprint/chap-Nagata-compactification-and-Chows-lemma/sec-Proof-of-Nagatas-theorem/_section.md), content hidden
\subsection*{\texorpdfstring{A.4\quad Proof of Nagata's theorem}{A.4 Proof of Nagata's theorem}}
\phantomsection\label{bp63}%
\addcontentsline{toc}{subsection}{\texorpdfstring{A.4\quad Proof of Nagata's theorem}{A.4 Proof of Nagata's theorem}}
% text: prose-A-23
\phantomsection\label{bp64}%
Over a Noetherian base, the theorem follows from Proposition \hyperref[bp61]{A.23} by induction over an affine cover whose members overlap densely. Each affine piece is compactified by its closure in a projective space. The proposition then compactifies the union of the first \(r\) pieces with the next one. It requires the overlap to be dense, because its valuative criterion only tests maps from spectra of valuation rings whose generic point maps into the overlap. Lemma \hyperref[bp65]{A.24} provides a cover whose total intersection contains all generic points of \(X\).

Over a general qcqs base we use Noetherian approximation. This works because a morphism of finite presentation is the base change of one over a Noetherian scheme, and compactifications are stable under base change. A morphism that is only of finite type, however, need not descend in this way. We therefore first embed \(X\) as a closed subscheme of a scheme of finite presentation over \(S\), and take the closure of \(X\) in a compactification of the latter.

% object: lemma-affine-covers-dense-intersection (blueprint/chap-Nagata-compactification-and-Chows-lemma/sec-Proof-of-Nagatas-theorem/lemma-affine-covers-dense-intersection/_section.md)
\phantomsection\label{bp65}%
\begin{bpstatementit}{Lemma~A.24}{ (Affine covers with dense intersection)}
For a scheme \(X\) of finite type over a Noetherian scheme \(S\), there is a finite affine cover \(X=\bigcup_{r=1}^m U_r\) whose total intersection is dense in \(X\). If \(X\to S\) is separated, then each \(U_r\to S\) has a projective compactification.
\end{bpstatementit}

% proof: proof/lemma-affine-covers-dense-intersection
\begin{proof}[{Proof}]\phantomsection\label{bp66}%
The scheme \(X\) has finitely many irreducible components. By [Stacks, \href{https://stacks.math.columbia.edu/tag/01ZX}{01ZX}], every point of \(X\) has an affine open neighbourhood that also contains all the generic points of those components, and since \(X\) is quasi-compact, finitely many such opens cover \(X\). Their total intersection still contains all generic points, so it is dense. For the second assertion, note that each \(U_r\) is affine of finite type over \(S\), so by [Stacks, \href{https://stacks.math.columbia.edu/tag/04II}{04II}] we may embed it as a locally closed subscheme of some \(\mathbf A^{n_r}_S\). Its schematic closure in \(\mathbf P^{n_r}_S\) is proper, and by Lemma \hyperref[bp9]{A.3} the scheme \(U_r\) is open in that closure.
\end{proof}

% text: prose-A-24
\phantomsection\label{bp67}%
We prove the theorem in three steps of increasing generality: a Noetherian base, a finitely presented morphism over any qcqs base, and a morphism of finite type.

% proof: proof/theorem-nagata-s-compactification-theorem
\begin{proof}[{Proof of Theorem A.2}]\phantomsection\label{bp68}%
If \(S\) is Noetherian, we choose a cover as in Lemma \hyperref[bp65]{A.24} and argue by induction on the unions \(X_r=U_1\cup\cdots\cup U_r\). For \(r=1\), the lemma gives a compactification of \(U_1\). If \(X_r\) has a compactification, we apply Proposition \hyperref[bp61]{A.23} to \(M=X_{r+1}\), with \(M_1=X_r\) and \(M_2=U_{r+1}\). Each piece has a compactification, by induction and by the lemma. The overlap \(X_r\cap U_{r+1}\) is dense in \(X_{r+1}\), since it contains the total intersection, which contains all generic points of \(X\). Hence \(X_{r+1}\) has a compactification. The last union is \(X\) itself, and taking the schematic closure of \(X\) makes the compactification dense.

Now let \(S\) be a general qcqs scheme, and suppose first that \(X\to S\) is separated and of finite presentation. By absolute Noetherian approximation we can write

\[
S=\varprojlim_i S_i
\]

with \(S_i\) Noetherian and affine transition maps [Stacks, \href{https://stacks.math.columbia.edu/tag/01ZA}{01ZA}]. Since \(X\to S\) is of finite presentation, it descends to a model \(X_i\to S_i\) for some \(i\), and after increasing \(i\) the model is separated [Stacks, \href{https://stacks.math.columbia.edu/tag/01ZM}{01ZM}, \href{https://stacks.math.columbia.edu/tag/01ZQ}{01ZQ}]. We then compactify this model by the Noetherian case and base change the compactification to \(S\). The result is a compactification, because open immersions and proper morphisms are stable under base change, and taking the schematic closure of \(X\) makes the result dense.

Finally, let \(X\to S\) be separated of finite type. By [Stacks, \href{https://stacks.math.columbia.edu/tag/01ZJ}{01ZJ}], there is a closed immersion \(X\hookrightarrow X_0\) with \(X_0\to S\) separated of finite presentation. We compactify \(X_0\) by the case already proved, and take the schematic closure of \(X\) in that compactification. This closure is proper over \(S\). Its restriction to the open \(X_0\) is the closed subscheme \(X\) itself, so \(X\) is open in the closure, and it is schematically dense there by Lemma \hyperref[bp9]{A.3}.
\end{proof}


% object: sec-Chows-lemma (blueprint/chap-Nagata-compactification-and-Chows-lemma/sec-Chows-lemma/_section.md), content hidden
\subsection*{\texorpdfstring{A.5\quad Chow's lemma}{A.5 Chow's lemma}}
\phantomsection\label{bp69}%
\addcontentsline{toc}{subsection}{\texorpdfstring{A.5\quad Chow's lemma}{A.5 Chow's lemma}}
% text: prose-A-25
\phantomsection\label{bp70}%
Let \(f:X\to S\) be proper. In Section \bpout{4.5} proper base change is first proved for projective morphisms. By Lemma \bpout{4.29} it then descends along a proper surjection \(\pi:\widetilde X\to X\) as soon as it holds for \(\pi\) and for the composite \(f\pi\). Chow's lemma says precisely that such a \(\pi\) exists with both \(\pi\) and \(f\pi\) projective. Some modification of \(X\) is unavoidable here, since there are proper varieties that are not projective.

% object: theorem-chow-s-lemma-projective-form (blueprint/chap-Nagata-compactification-and-Chows-lemma/sec-Chows-lemma/theorem-chow-s-lemma-projective-form/_section.md)
\phantomsection\label{bp71}%
\begin{bpstatementit}{Theorem~A.25}{ (Chow's lemma, projective form)}
Let \(f:X\to S\) be proper with \(S\) qcqs. There are a proper surjection \(\pi:\widetilde X\to X\) and a closed immersion \(i_S:\widetilde X\hookrightarrow\mathbf P^n_S\) such that \(f\pi=q_Si_S\), where \(q_S:\mathbf P^n_S\to S\) is the projection. Moreover,

\[
i_X=(\pi,i_S):\widetilde X\hookrightarrow
X\times_S\mathbf P^n_S=\mathbf P^n_X
\]

is a closed immersion. Both \(\pi\) and \(f\pi\) are projective.
\end{bpstatementit}

% proof: proof/theorem-chow-s-lemma-projective-form
\begin{proof}[{Proof}]\phantomsection\label{bp72}%
We take the immersion into projective space given by the Stacks Project and show that properness of \(X\) makes it closed. The second assertion then follows by writing \(i_X\) as a graph followed by a base change of \(i_S\).

The version of Chow's lemma in [Stacks, \href{https://stacks.math.columbia.edu/tag/0202}{0202}] gives a proper surjection \(\pi:\widetilde X\to X\) and an immersion \(i_S:\widetilde X\to\mathbf P^n_S\) for every separated \(X\to S\) of finite type with \(S\) qcqs. It is deduced from the Noetherian case by the approximation argument in the proof of Theorem \hyperref[bp7]{A.2}, using the embedding of \(X\) into a scheme of finite presentation [Stacks, \href{https://stacks.math.columbia.edu/tag/01ZJ}{01ZJ}]. In our situation, moreover, \(f\) and \(\pi\) are proper, so \(\widetilde X\) is proper over \(S\). Now a map from a proper \(S\)-scheme to a separated \(S\)-scheme is proper, because its graph is closed and the projection of the product to the target is proper. Hence \(i_S\) is a proper immersion. Its image is therefore closed, and \(i_S\) is a closed immersion.

For the second assertion, we factor \(i_X\) as the graph of \(\pi\) followed by the base change of \(i_S\):

\[
\widetilde X\longrightarrow\widetilde X\times_SX
\longrightarrow\mathbf P^n_S\times_SX.
\]

The first arrow is closed because \(X\to S\) is separated, and the second is closed because \(i_S\) is. Their composite is \(i_X\). Since \(\pi\) is \(i_X\) followed by the projection \(\mathbf P^n_X\to X\), it is projective, and so is \(f\pi=q_Si_S\).
\end{proof}



\clearpage\phantomsection\hypertarget{bp-outside}{}%
\section*{Outside this export}
\addcontentsline{toc}{section}{Outside this export}
References marked \ensuremath{{}^\ast} are to parts of the blueprint this document leaves out:
\begin{description}
\item[Proposition 2.58] Extension by zero and restriction to a closed subscheme, in Ⅰ The six operations (\bpslug{proposition-{}\allowbreak{}extension-{}\allowbreak{}zero-{}\allowbreak{}restriction-{}\allowbreak{}closed-{}\allowbreak{}subscheme})
\item[Section 4.5] Proper morphisms, in Ⅰ The six operations (\bpslug{sec-{}\allowbreak{}Proper-{}\allowbreak{}morphisms})
\item[Lemma 4.29] Composition and cancellation, in Ⅰ The six operations (\bpslug{lemma-{}\allowbreak{}composition-{}\allowbreak{}cancellation})
\item[Proposition 5.26] Construction and independence of $Rf_!$, in Ⅰ The six operations (\bpslug{proposition-{}\allowbreak{}construction-{}\allowbreak{}independence-{}\allowbreak{}rf-{}\allowbreak{}shriek})
\end{description}
\clearpage\phantomsection%
\section*{References}
\addcontentsline{toc}{section}{References}
\bpbibitem{BBD82}{bp74}{A. Beilinson, J. Bernstein, and P. Deligne. \emph{Faisceaux pervers}, volume 100 of \emph{Astérisque}. 1982.}

\bpbibitem{BS15}{bp75}{B. Bhatt and P. Scholze. The pro-étale topology for schemes. \emph{Astérisque}, 369:99–201, 2015. \href{https://people.mpim-bonn.mpg.de/scholze/proetale.pdf}{Author's manuscript}.}

\bpbibitem{Con07}{bp73}{B. Conrad. Deligne's notes on Nagata compactifications. \emph{Journal of the Ramanujan Mathematical Society}, 22:205–257, 2007. \href{https://math.stanford.edu/\~conrad/papers/nagatafinal.pdf}{Author's manuscript}.}

\bpbibitem{Del74}{bp76}{P. Deligne. La conjecture de Weil. I. \emph{Publ. Math. IHÉS}, 43:273–307, 1974.}

\bpbibitem{Del80}{bp77}{P. Deligne. La conjecture de Weil. II. \emph{Publ. Math. IHÉS}, 52:137–252, 1980.}

\bpbibitem{Den22}{bp78}{C. Deninger. There is no “Weil-”cohomology theory with real coefficients for arithmetic curves. \emph{Ann. Sc. Norm. Super. Pisa Cl. Sci.}, 2022. To appear. arXiv:\href{https://arxiv.org/abs/2204.02714}{2204.02714}.}

\bpbibitem{Deu41}{bp79}{M. Deuring. Die Typen der Multiplikatorenringe elliptischer Funktionenkörper. \emph{Abh. Math. Sem. Univ. Hamburg}, 14:197–272, 1941.}

\bpbibitem{Eke90}{bp80}{T. Ekedahl. On the adic formalism. In \emph{The Grothendieck Festschrift, Vol. II}, volume 87 of \emph{Progress in Mathematics}, pages 197–218. Birkhäuser, 1990.}

\bpbibitem{Has36}{bp81}{H. Hasse. Zur Theorie der abstrakten elliptischen Funktionenkörper. I, II, III. \emph{J. Reine Angew. Math.}, 175:55–62, 69–88, 193–208, 1936.}

\bpbibitem{Hum72}{bp82}{J. E. Humphreys. \emph{Introduction to Lie Algebras and Representation Theory}, volume 9 of \emph{Graduate Texts in Mathematics}. Springer, 1972.}

\bpbibitem{Ill94}{bp83}{L. Illusie. Autour du théorème de monodromie locale. \emph{Astérisque}, 223:9–57, 1994. §4. \href{https://www.numdam.org/item/AST_1994__223__9_0/}{Original article}.}

\bpbibitem{Kat88}{bp84}{N. M. Katz. Travaux de Laumon. \emph{Astérisque}, 161–162:105–132, 1988. Séminaire Bourbaki, exp. 691.}

\bpbibitem{KW01}{bp85}{R. Kiehl and R. Weissauer. \emph{Weil conjectures, perverse sheaves and \(\ell\)-adic Fourier transform}, volume 42 of \emph{Ergebnisse der Mathematik}. Springer, 2001.}

\bpbibitem{Lau87}{bp86}{G. Laumon. Transformation de Fourier, constantes d'équations fonctionnelles et conjecture de Weil. \emph{Publ. Math. IHÉS}, 65:131–210, 1987.}

\bpbibitem{Mil86a}{bp87}{J. S. Milne. \href{https://www.jmilne.org/math/xnotes/AVs.pdf}{Abelian Varieties}. In \emph{Arithmetic Geometry}. 1986. Corrected online version.}

\bpbibitem{Mil86b}{bp88}{J. S. Milne. \href{https://www.jmilne.org/math/xnotes/JVs.pdf}{Jacobian Varieties}. In \emph{Arithmetic Geometry}. 1986. Corrected online version.}

\bpbibitem{Ran39}{bp89}{R. A. Rankin. Contributions to the theory of Ramanujan's function \(\tau(n)\) and similar arithmetical functions. II. \emph{Proc. Cambridge Philos. Soc.}, 35:357–372, 1939.}

\bpbibitem{SGA4}{bp90}{M. Artin, A. Grothendieck, and J.-L. Verdier, editors. \emph{Théorie des topos et cohomologie étale des schémas}, volumes 269, 270, 305 of \emph{Lecture Notes in Mathematics}. Springer, 1972–1973. SGA 4.}

\bpbibitem{SGA4½}{bp91}{P. Deligne. \emph{Cohomologie étale}, volume 569 of \emph{Lecture Notes in Mathematics}. Springer, 1977. SGA 4½; \emph{Théorèmes de finitude}, §3, and \emph{Rapport sur la formule des traces}, §3, and \emph{Sommes trigonométriques}, §1.}

\bpbibitem{SGA5}{bp92}{A. Grothendieck, editor. \emph{Cohomologie \(\ell\)-adique et fonctions \(L\)}, volume 589 of \emph{Lecture Notes in Mathematics}. Springer, 1977. SGA 5.}

\bpbibitem{SGA7II}{bp93}{P. Deligne and N. Katz, editors. \emph{Groupes de monodromie en géométrie algébrique}, volume 340 of \emph{Lecture Notes in Mathematics}. Springer, 1973. SGA 7 II; Exposé XIII, (2.1.7).}

\bpbibitem{Stacks}{bp94}{The Stacks Project Authors. \href{https://stacks.math.columbia.edu}{The Stacks Project}. Individual results are cited by their stable tags in the text.}

\bpbibitem{Var05}{bp95}{Y. Varshavsky. Lefschetz–Verdier trace formula and a generalization of a theorem of Fujiwara. 2005. Version 2. arXiv:\href{https://arxiv.org/abs/math/0505564}{math/0505564}. \href{https://arxiv.org/abs/math/0505564v2}{Author's preprint}.}

\bpbibitem{Wei48a}{bp96}{A. Weil. \emph{Sur les courbes algébriques et les variétés qui s'en déduisent}. Hermann, Paris, 1948.}

\bpbibitem{Wei48b}{bp97}{A. Weil. \emph{Variétés abéliennes et courbes algébriques}. Hermann, Paris, 1948.}

\bpbibitem{Wei49}{bp98}{A. Weil. Numbers of solutions of equations in finite fields. \emph{Bull. Amer. Math. Soc.}, 55:497–508, 1949.}

\end{document}
